% source: 490b60895ca6eb6499ecf5a0faf0563b02e5fb47
\chapter{格}
\label{ch:lattice}

\section{本章的意义}
本章是整门课程中最重要的章节之一。如果说前面几章介绍了向量、矩阵和多项式环，那么本章要解释的是使后量子密码成为可能的那个几何对象：格（lattice）。格是 LWE、RLWE、Kyber、Dilithium 以及许多其他现代构造背后的核心结构。

\section{学习目标}
读完本章后，你应当能够：

\begin{itemize}
  \item 解释什么是格；
  \item 区分格与向量空间；
  \item 理解为什么一个格可以有许多不同的基；
  \item 理解好基与坏基之间的区别；
  \item 体会为什么 SVP 和 CVP 这类问题是困难的；
  \item 看出格问题与 LWE 之间的联系。
\end{itemize}

\section{从二维开始}
我们不直接跳入高维理论，而是从一个简单的二维例子开始。设

\[
b_1 = \begin{pmatrix}2\\1\end{pmatrix}, \qquad b_2 = \begin{pmatrix}1\\3\end{pmatrix}.
\]

现在只允许整数系数：

\[
v = x b_1 + y b_2, \qquad x,y \in \mathbb{Z}.
\]

如果取许多这样的组合，所得的点会在平面上构成一个由离散点组成的无限规则图案。这个集合称为格。

\begin{figure}[H]
\centering
\begin{tikzpicture}[scale=0.45]
  \fill[pqclight] (0,0) -- (2,1) -- (3,4) -- (1,3) -- cycle;
  \draw[pqcgray!40, thin] (-6.5,0) -- (6.5,0);
  \draw[pqcgray!40, thin] (0,-8.5) -- (0,8.5);
  \foreach \i in {-2,...,2} {
    \foreach \j in {-2,...,2} {
      \pgfmathsetmacro{\x}{2*\i+1*\j}
      \pgfmathsetmacro{\y}{1*\i+3*\j}
      \node[latticept] at (\x,\y) {};
    }
  }
  \node[latticeptbold] at (0,0) {};
  \draw[basisvec] (0,0) -- (2,1) node[midway, below right, font=\scriptsize] {$b_1$};
  \draw[altvec] (0,0) -- (1,3) node[midway, above left, font=\scriptsize] {$b_2$};
\end{tikzpicture}
\caption{由 $b_1=(2,1)$ 与 $b_2=(1,3)$ 的所有整数组合 $xb_1+yb_2$ 生成的格 $L(b_1,b_2)$。阴影平行四边形是由该基张成的基本区域。}
\label{fig:lattice-basic}
\end{figure}

\subsection{Sage 实验：生成一个二维格}

\begin{lstlisting}[language=Python]
b1 = vector(ZZ, [2, 1])
b2 = vector(ZZ, [1, 3])
points = []
for i in range(-5, 6):
    for j in range(-5, 6):
        points.append(i * b1 + j * b2)
points[:10]
\end{lstlisting}

输出是一个格点列表。重要的观察是：这些格点以非常规则、离散的方式排列。

\section{形式化定义}
形式地说，若 $b_1, \ldots, b_n$ 是线性无关的向量，则集合

\[
L(B) = \left\{ \sum_{i=1}^{n} z_i b_i \;:\; z_i \in \mathbb{Z} \right\}
\]

称为由基 $B = \{b_1, \ldots, b_n\}$ 生成的格。

与向量空间的关键区别在于：系数被限制为整数，而不是任意实数。

这个定义可以直接执行：在有界区域内枚举一个格，恰好就是对整数系数向量做有界搜索。

\begin{algorithm}[H]
\caption{有界格点枚举}
\label{alg:lattice-enum}
\begin{algorithmic}[1]
\Require 基 $b_1,\dots,b_n \in \mathbb{Z}^n$；系数界 $N \in \mathbb{Z}_{>0}$
\Ensure 所有满足 $|z_i| \le N$ 的格点 $\sum_i z_i b_i$ 构成的集合 $S$
\State $S \gets \emptyset$
\For{each $(z_1,\dots,z_n) \in \{-N,\dots,N\}^n$}
    \State $v \gets \sum_{i=1}^{n} z_i \, b_i$
    \State $S \gets S \cup \{v\}$
\EndFor
\State \Return $S$
\end{algorithmic}
\end{algorithm}

这正是上面的 Sage 实验在 $n=2$、$N=5$ 时所做的：对整数系数做两重嵌套循环，收集所得的每一个点。图~\ref{fig:lattice-basic} 中的规则网格就是这个算法的输出。

\section{为什么格不是向量空间}
向量空间允许任意的实系数或复系数。例如，

\[
0.13 b_1 + 2.78 b_2
\]

是由 $b_1$ 与 $b_2$ 张成的向量空间中一个完全合法的元素。但在格中，系数必须是整数。这就是为什么格是离散结构而非连续结构。

这一区别至关重要。密码学需要离散的问题，因为连续的问题往往很容易求解。

\section{向量空间与格的对比}
向量空间是连续而稠密的。格是离散而稀疏的。换言之：

\begin{itemize}
  \item 向量空间：所有实系数的线性组合；
  \item 格：仅整数线性组合。
\end{itemize}

正是这一差别使格问题既困难又可用于密码学。

\section{一个格可以有许多基}
最重要的概念要点是：一个格可以用许多不同的基来描述。例如，基

\[
\left\{(1,0), (0,1)\right\}
\]

与基

\[
\left\{(2,1), (1,1)\right\}
\]

可以生成同一个格。

这就好比一座城市可以用不同的坐标系来描述。城市本身没有变，变的只是坐标描述。

\begin{figure}[H]
\centering
\begin{tikzpicture}[scale=0.75]
  \draw[pqcgray!40, thin] (-3.5,0) -- (3.5,0);
  \draw[pqcgray!40, thin] (0,-3.5) -- (0,3.5);
  \foreach \i in {-3,...,3} {
    \foreach \j in {-3,...,3} {
      \node[latticept] at (\i,\j) {};
    }
  }
  \node[latticeptbold] at (0,0) {};
  \draw[basisvec] (0,0) -- (1,0);
  \node[pqcblue, font=\small] at (1.0,-0.35) {$e_1$};
  \draw[basisvec] (0,0) -- (0,1);
  \node[pqcblue, font=\small] at (-0.35,1.0) {$e_2$};
  \draw[altvec] (0,0) -- (2,1);
  \node[pqcteal, font=\small] at (2.25,0.75) {$b_1'$};
  \draw[altvec] (0,0) -- (1,1);
  \node[pqcteal, font=\small] at (0.65,1.35) {$b_2'$};
\end{tikzpicture}

{\small\centering $e_1=(1,0),\ e_2=(0,1) \qquad b_1'=(2,1),\ b_2'=(1,1)$\par}
\vspace{2pt}
\caption{同一个格 $\mathbb{Z}^2$ 的两组不同的基：蓝色的标准基 $\{e_1,e_2\}$ 与青色的倾斜基 $\{b_1',b_2'\}$。由于 $\det\left(\begin{smallmatrix}2&1\\1&1\end{smallmatrix}\right)=1$，两者都能到达同一底层网格中的每一个点。}
\label{fig:lattice-two-bases}
\end{figure}

\subsection{Sage 实验：两组不同的基}

\begin{lstlisting}[language=Python]
B1 = matrix(ZZ, [[1, 0], [0, 1]])
B2 = matrix(ZZ, [[2, 1], [1, 1]])
\end{lstlisting}

稍后我们将验证它们生成的是同一个格。

\section{好基与坏基}
现在我们来到格基约化背后的核心思想。有些基是好的，有些基是坏的。

坏基由近乎平行且很长的向量组成。这样的基使人很难找到邻近的格点。

好基由相对较短且近乎正交的向量组成，使人更容易在格中“导航”。

这就是 LLL 和 BKZ 等算法存在的原因：它们试图把一组坏基变换为一组更好的基。

\begin{figure}[H]
\centering
\begin{subfigure}[t]{0.47\textwidth}
\centering
\begin{tikzpicture}[scale=0.30]
  \draw[pqcgray!40, thin] (-7.5,0) -- (7.5,0);
  \draw[pqcgray!40, thin] (0,-7.5) -- (0,7.5);
  \foreach \i in {-7,...,7} {
    \foreach \j in {-7,...,7} {
      \node[latticept] at (\i,\j) {};
    }
  }
  \node[latticeptbold] at (0,0) {};
  \draw[shortvec] (0,0) -- (1,0);
  \node[pqcorange, font=\scriptsize] at (1.6,-0.5) {$(1,0)$};
  \draw[shortvec] (0,0) -- (0,1);
  \node[pqcorange, font=\scriptsize] at (-0.9,1.5) {$(0,1)$};
\end{tikzpicture}
\caption{好基：短且近乎正交。}
\end{subfigure}
\hfill
\begin{subfigure}[t]{0.47\textwidth}
\centering
\begin{tikzpicture}[scale=0.30]
  \draw[pqcgray!40, thin] (-7.5,0) -- (7.5,0);
  \draw[pqcgray!40, thin] (0,-7.5) -- (0,7.5);
  \foreach \i in {-7,...,7} {
    \foreach \j in {-7,...,7} {
      \node[latticept] at (\i,\j) {};
    }
  }
  \node[latticeptbold] at (0,0) {};
  \draw[altvec] (0,0) -- (4,3);
  \node[pqcteal, font=\scriptsize] at (4.9,2.6) {$(4,3)$};
  \draw[altvec] (0,0) -- (7,5);
  \node[pqcteal, font=\scriptsize] at (6.1,5.7) {$(7,5)$};
\end{tikzpicture}
\caption{\emph{同一个}格 $\mathbb{Z}^2$ 的坏基：长且近乎平行。}
\end{subfigure}
\caption{两组基生成的是完全相同的格（两种情形下均有 $\det=\pm1$），但右侧的坏基几乎让格的几何结构无从看出——两个向量几乎指向同一方向。从右侧这样的坏基恢复出左侧那样的好基，恰恰就是格基约化所要解决的问题。}
\label{fig:good-bad-basis}
\end{figure}

\section{基本平行多面体}
给定一组基，就可以构成基本平行多面体，即格的基本单元。它的体积与基矩阵的行列式有关。

对于二维格，这个单元的面积由行列式的绝对值给出。

\subsection{Sage 实验：行列式}

\begin{lstlisting}[language=Python]
B = matrix(ZZ, [[2, 1], [1, 3]])
B.det()
\end{lstlisting}

行列式衡量的是格的稠密程度。行列式越大，通常意味着格越稀疏；行列式越小，则意味着格越稠密。

\section{格的维数}
在二维中，格很容易可视化。然而在密码学中，维数要大得多。Kyber、Dilithium 及其他基于格的系统可能在数百乃至数千维中工作。在这样的高维空间中，几何图景要难以想象得多，但许多相同的直觉依然适用。

\section{为什么格问题会变得困难}
在低维中，人们常常可以手工或借助简单的几何论证来求解格问题。然而在更高维中，寻找短向量或邻近向量会变得极其困难。这种复杂度的爆炸式增长正是基于格的安全性的基础。

\section{Sage 实验：最短向量}
Sage 提供了内置的格工具。可以按如下方式直接由整数矩阵创建一个格对象：

\begin{lstlisting}[language=Python]
from sage.modules.free_module_integer import IntegerLattice

B = matrix(ZZ, [[2, 1], [1, 3]])
L = IntegerLattice(B)
L
\end{lstlisting}

定义好格之后，只需一条命令即可计算最短向量：

\begin{lstlisting}[language=Python]
L.shortest_vector()
\end{lstlisting}

这就是最短向量问题（SVP）背后的基本对象，而 SVP 是基于格的密码学中核心的困难性假设之一。

\section{实验}

ML-KEM（FIPS 203~\cite{fips203}）与 ML-DSA（FIPS 204~\cite{fips204}）底层的格位于数百维的空间中，其几何结构无法可视化。下面的实验在二维中——在那里每个对象都可以画出来并手工核验——分离出这些标准最终所依赖的四个现象：点集的离散性、行列式的基不变性、好基与坏基之间的差距，以及最短向量的难解性。每一个现象在原理上都原封不动地出现在密码学规模中。

\subsection{实验 1：枚举格}
生成格 $\Lambda(\mathbf{b}_1,\mathbf{b}_2)$ 的一个有限窗口，并确认它是一个离散的点阵网格，而不是连续的张成空间。

\begin{lstlisting}[language=Python]
b1 = vector(ZZ, [2, 1])
b2 = vector(ZZ, [1, 3])
pts = [i * b1 + j * b2 for i in range(-4, 5) for j in range(-4, 5)]
pts
\end{lstlisting}

绘制 \texttt{pts} 即可看到典型的规则网格。正是这同一个对象，在维数为 $256$ 且定义在 $\mathbb{Z}_q$ 而非 $\mathbb{Z}$ 上时，成为 ML-KEM 核心的模格。

\subsection{实验 2：行列式是基不变量}
对\emph{同一个}格的若干组基计算格行列式 $\det\Lambda = |\det B|$，每组基都由 $B$ 左乘一个幺模矩阵 $U \in \mathrm{GL}_n(\mathbb{Z})$（即满足 $\det U = \pm 1$ 的整数矩阵）得到。

\begin{lstlisting}[language=Python]
B = matrix(ZZ, [[2, 1], [1, 3]])
U = matrix(ZZ, [[1, 0], [1, 1]])          # unimodular: det U = 1
print(abs(B.det()), abs((U * B).det()))   # identical
\end{lstlisting}

行列式——即基本平行多面体的余体积，因而也是格密度的倒数——保持不变，而基向量本身则完全改变。这种在 $\mathrm{GL}_n(\mathbb{Z})$ 作用下的不变性，恰恰使一个方案能够公开格的一组基，而秘密持有另一组基。

\subsection{实验 3：好基与坏基}
同一个格既有近乎正交的（“好”）基，也有高度倾斜的（“坏”）基。对比两者，并通过正交性缺陷 $\prod_i \lVert \mathbf{b}_i\rVert / \det\Lambda$ 来量化倾斜程度。

\begin{lstlisting}[language=Python]
good = matrix(ZZ, [[1, 0], [0, 1]])       # orthonormal basis of Z^2
bad  = matrix(ZZ, [[1, 7], [0, 1]])       # a sheared basis of the same Z^2
\end{lstlisting}

这一差距是陷门格密码的核心。在基于 NTRU 的 FN-DSA 中，以及在精神上贯穿整个 ML-DSA（FIPS 204~\cite{fips204}）的设计中，秘密钥充当好基，公钥充当坏基；在没有好基的情况下伪造签名，可归约为一个困难的格问题。第~\ref{ch:lll}~章将把好基与坏基的区别定量化，并说明一组坏基能被约化到多么接近好基的程度。

\subsection{实验 4：最短向量是格不变量}
对同一个格的若干组不同的基计算最短非零向量，并确认其长度是格的性质，而不是所选基的性质。

\begin{lstlisting}[language=Python]
from sage.modules.free_module_integer import IntegerLattice
B1 = matrix(ZZ, [[2, 1], [1, 3]])
B2 = matrix(ZZ, [[3, 4], [4, 7]])          # = U*B1, U unimodular: same lattice
for B in [B1, B2]:
    print(IntegerLattice(B).shortest_vector())   # identical length
\end{lstlisting}

在二维中，\texttt{shortest\_vector} 瞬间就能给出答案。而在 FIPS 203 与 204 所涉及的维数下，没有任何已知算法能在可行时间内做到这一点；这种渐近意义上的难解性——形式化为最短向量问题（SVP）及其近似版本和判定版本——正是这两项标准共同依托的安全基础。

\section{小结}
本章的主要结论有三点：

\begin{enumerate}
  \item 格不是向量空间；它是由整数线性组合构成的离散结构。
  \item 一个格可以有许多不同的基，而基的选择会极大地影响处理这个格的难易程度。
  \item 格问题（尤其是 SVP 和 CVP）的困难性是许多后量子密码系统的基础。
\end{enumerate}

下一章\emph{Gram--Schmidt 正交化与 LLL 算法}将介绍第一个把坏基变为好基的主要算法工具。

\printbibliography[heading=chapterrefs]
